\documentclass[a4paper,12pt]{article}
\usepackage{listings}
\usepackage{xcolor}
\usepackage{hyperref}
\usepackage{geometry}
\geometry{margin=1in}

\title{JavaScript: An Overview}
\author{}
\date{March 5, 2025}

\begin{document}

\maketitle

\section{Introduction}
JavaScript is responsible for the interactive layer of a web page, complementing the structural layer provided by markup and the presentational layer provided by CSS. JavaScript allows developers to modify the structure and presentation of a page by adding, removing, and altering markup and styles in response to user interaction and internal logic.

\section{History of JavaScript}
JavaScript's syntax was initially inspired by Java's syntax. It was originally named \textbf{LiveScript} when it first shipped in a beta version of Netscape Navigator in 1995. Later, it was renamed \textbf{JavaScript}. The current LiveScript is a language that generates JavaScript.

JavaScript was created in 1995 by \textbf{Brendan Eich}, a software engineer at Netscape Communications Corporation. Microsoft released their own implementation of JavaScript, called \textbf{JScript}, shortly afterward with Internet Explorer 3.0.

\begin{itemize}
    \item \textbf{Creator}: Brendan Eich
    \item \textbf{Trademark Holder}: Oracle Corporation
    \item \textbf{Originating Company}: Sun Microsystems
    \item \textbf{Current Version (as of March 5, 2025)}: ECMAScript 2023 (ES14), released in June 2023.
\end{itemize}

ECMAScript is the standard that informs the creation of specific scripting languages such as Microsoft's JScript, Adobe's ActionScript, and JavaScript itself.

\section{Execution Model}
JavaScript is executed using a combination of \textbf{interpretation} and \textbf{Just-In-Time (JIT) compilation}. This hybrid approach provides the flexibility of interpretation with the performance benefits of compilation.

A script is sent to the browser alongside assets such as markup, images, and stylesheets. The browser interprets it as a human-readable sequence of Unicode characters, parsed from left to right and top to bottom. The JavaScript interpreter performs lexical analysis, parsing the script into discrete input elements:

\begin{itemize}
    \item Tokens
    \item Format control characters
    \item Line terminators
    \item Comments
    \item Whitespace
\end{itemize}

JavaScript changes disappear when you reload or move to a new page unless explicitly saved.

\section{JavaScript Fundamentals}
\subsection{Statements}
A statement is a unit of instruction made up of one or more lines of code representing an action.

\begin{lstlisting}[language=JavaScript, caption=Example of a Statement]
let myVariable = 4;
myVariable;
\end{lstlisting}

To be interpreted correctly, statements must end in a semicolon. However, JavaScript includes a feature called \textbf{automatic semicolon insertion}, which treats a line break as a semicolon if a missing semicolon would cause an error.

\subsection{Block Statements}
A block statement groups multiple statements inside curly braces (\{\}). It allows combining statements where JavaScript expects only one.

\begin{lstlisting}[language=JavaScript, caption=Example of a Block Statement]
if (x === 2) {
  // Some behavior;
}
\end{lstlisting}

\subsection{Expressions}
An expression is a unit of code that results in a value.

\begin{lstlisting}[language=JavaScript, caption=Example of an Expression]
2 + 2; // Results in 4
\end{lstlisting}

\subsubsection{Grouping Operator}
Parentheses \texttt{()} are used to group parts of an expression.

\begin{lstlisting}[language=JavaScript, caption=Grouping Operator Example]
2 + 2 * 4; // Results in 10
(2 + 2) * 4; // Results in 16
\end{lstlisting}

\subsection{Type Coercion}
JavaScript is a weakly typed language, meaning data values do not need to be explicitly marked with a specific data type.

\begin{lstlisting}[language=JavaScript, caption=Example of Weak Typing]
"1" + 1; // Results in "11"
\end{lstlisting}

Datatypes can also be explicitly coerced using JavaScript's built-in methods.

\section{JavaScript Syntax}
\subsection{Case Sensitivity}
JavaScript is fully case-sensitive.

\begin{lstlisting}[language=JavaScript, caption=Case Sensitivity Example]
console.log("Log this."); // Works
console.Log("Log this too."); // Throws an error
\end{lstlisting}

\subsection{Whitespace Insensitivity}
JavaScript ignores the amount and type of whitespace.

\begin{lstlisting}[language=JavaScript, caption=Whitespace Insensitivity Example]
console.log("Log this"); console.log("Log this too");
\end{lstlisting}

However, whitespace can be significant as a separator between lexical tokens.

\begin{lstlisting}[language=JavaScript, caption=Whitespace as a Separator]
let x;  // Valid
letx;   // Error: letx is not defined
\end{lstlisting}

Whitespace is often used to improve code readability through indentation.

\end{document}
